\section{Calibration}\label{sec:arch-calibration-chapter}
How the models are fitted to the market, how the calibrated models are stored and reused, and how often each part is
refreshed. The equity vol surfaces themselves (from listed options) are market data: Section~\ref{sec:arch-equity}.

% ─────────────────────────────────────────────────────────────────────────────
\subsection{Two ways to calibrate}\label{sec:arch-calibration}
Every equity model can be used in two ways.
\begin{itemize}[nosep]
  \item \emph{Inside the pricing} (\code{Model.from\_equity(equity, maturity, strike)}, a model class given to a
        pricer): the calibration runs at every pricing call, inside the pricer's tape. The price always uses today's
        parameters and its market greeks include the move of the parameters (greeks through the calibration). Simple,
        but the calibration time is paid at every price.
  \item \emph{Calibrated once and held} (\code{Model.calibrate(equity)} $\to$ an \code{EquityModel} instance given to a
        pricer): the parameters are fixed numbers (\code{tf.Variable}s) bound to today's market by
        \code{model.process(equity)}. Pricing is fast; the greeks are model greeks (\code{model.parameters}) plus the
        direct market greeks, and \code{par\_greeks} maps the model greeks to the market with the calibration Jacobian
        stored at calibration time.
\end{itemize}
For the models of the calibration sets, the two give the same greeks on the calibration market
(\code{tests/test\_model\_instances.py}); the instance diverges from the market as the market moves away from the
calibration date, which is measured by \code{check\_calibration}.

\paragraph{Differentiable calibration (implicit function).} The least-squares calibrations (Heston, displaced diffusion) run in NumPy / SciPy and end with the implicit step
(\code{\_implicit\_step}, helpers in \code{numericalhandles/implicit.py}): at the optimum the gradient $g$ of the
objective, its Hessian $H$ and $\partial g/\partial m$ for the calibration inputs $m$ (relative strikes, target prices,
weights) are computed once, eagerly, as numbers (\code{optimality\_derivatives}: one backward pass per parameter);
\code{implicit\_solution} returns the parameters as a tensor of $m$ with a custom gradient
$\mathrm{d}p/\mathrm{d}m = -H^{-1}\,\partial g/\partial m$ (parameters on a bound fixed). The caller's tape never
differentiates through the second-order graph of $g$ (first order only). Inside \code{from\_equity} this makes the
parameters differentiable functions of the market; \code{calibrate} takes one Jacobian per parameter to all the risk
factors and stores it (\code{EquityModel.jacobian}), so that the par greeks of an instance equal the greeks through the
calibration (\code{tests/test\_model\_instances.py}).

% ─────────────────────────────────────────────────────────────────────────────
\subsection{What each equity model fits}\label{sec:arch-equity-calibration}
\begin{center}
\small
\begin{tabular}{@{}L{3.1cm}L{12.1cm}@{}}
\toprule
Model & Calibration \\
\midrule
Local vol & none: the local vol is computed from the surface at every pricing (the model has only settings) \\
Displaced diffusion & per surface maturity, $(\beta, \sigma)$ by least squares on the call prices per unit forward
  (bounded; \code{\_fit\_slices}), then the implicit step and the forward pieces (\code{forward\_from\_slices}) \\
Heston & the five parameters on the whole surface (strikes per unit forward, call prices, each maturity weighted by
  $1/$ATM vega), bounded least squares on a cheap COS then polished on the full-precision COS, then the implicit
  step; start from the ATM vols ($\kappa = 1.5$, $\xi = 0.5$, $\rho = -0.6$) or \code{initial} \\
Heston SLV & \code{HestonModel.calibrate} (Heston part), the local vol settings, then \code{calibrate\_leverage}: the
  $\E[v \mid X]$ table (50\,000 particles, 50 rows a year to the last surface maturity, 41 smoother nodes per row)
  and, with \code{bias\_correction}, the Richardson-corrected table that gives the price \\
\bottomrule
\end{tabular}
\end{center}
The fit is reported by \code{model\_fit(model, equity)} (model and market vols on the surface nodes: closed forms for
Heston and displaced diffusion, Monte Carlo for Heston SLV). Quality is the rms error in vol points where a diffusive
model is meaningful (\code{\_fit\_errors}: maturities from one month, strikes within two standard deviations): good up
to 1, fair up to 2.5, poor above. Typical timings on the real snapshot: displaced diffusion 30--40~s, Heston a few
seconds, Heston SLV about 20~s.

% ─────────────────────────────────────────────────────────────────────────────
\subsection{Calibration sets}\label{sec:arch-calibration-sets}
The models are calibrated less often than the market is downloaded, and stored apart from the snapshots
(\code{calibration/models.py}):
\begin{itemize}[nosep]
  \item \code{calibrate\_snapshot\_models(snapshot, root, plan=...)} calibrates the equities of a snapshot, one model
        at a time (a failure is recorded and does not stop the set), and writes
        \code{<root>/<yyyymmdd>/<hhmm>/}: \code{calibration\_info.json} (source snapshot, models, times),
        \code{summary.csv} (per underlying and model: status ok / error / skipped, quality, fit errors, seconds,
        parameters), and per underlying \code{<ticker>/<model>\_<table>.csv} (the tables of \code{to\_frames}: settings,
        parameters, Jacobian, report; for Heston SLV also the nested Heston and local vol tables and the leverage
        nodes) and \code{<ticker>/<model>\_fit.csv};
  \item \code{list\_calibration\_sets}, \code{load\_calibration\_set}, \code{load\_model(folder, ticker, name)}: read
        them back (\code{from\_frames}, exact round trip);
  \item \code{check\_calibration(model, equity)}: the fit of stored parameters on a newer market (\code{stale} above
        \code{STALE\_VOLPTS} = 1 vol point);
  \item \code{model\_test(models, equity, mkt)}: vanillas, down-and-in puts observed at expiry and cash digitals,
        priced with each model on shared paths, against Black and local vol.
\end{itemize}
The sets live in \code{marketdata/calibrations/<source>/} (git-ignored). They are run from the command line
(\code{python -m tqstudio.calibrate}, options \code{--date}, \code{--set}, \code{--models}, \code{--tickers},
\code{--slv}) or from the tqstudio Calibration tab. Which models are calibrated for which underlying is the
\emph{universe}: \code{marketdata/config/calibration.yaml}, one flag per model and underlying by asset class
(defaults per asset class; edited in the Universe sub-tab; rates models listed but off), turned into the
\code{plan} \code{\{ticker: [models]\}}. The tqstudio Tradepad prices with the models of the latest set that holds the
underlying (Heston, displaced diffusion, Heston SLV), bound to the current market.

% ─────────────────────────────────────────────────────────────────────────────
\subsection{Refresh frequency: what moves when}\label{sec:arch-calibration-frequency}
A model instance is several layers that need not be refreshed together:
\begin{center}
\small
\begin{tabular}{@{}L{4.2cm}L{4.2cm}L{6.8cm}@{}}
\toprule
Layer & Refreshed & How \\
\midrule
Forward (spot, repo, dividends, curve) & every pricing & read from the market by \code{process(equity)} \\
Local vol (and the local vol part of SLV) & every pricing & rebuilt from today's surface by \code{process} \\
Heston, displaced diffusion parameters & daily, with a calibration set & \code{calibrate}; held between sets \\
SLV table $\E[v \mid X]$ & when the surface has moved & \code{calibrate\_leverage(equity)}, Heston held \\
SLV Heston part & rarely (weekly, or on a trigger) & \code{HestonModel.calibrate} \\
\bottomrule
\end{tabular}
\end{center}
Heston SLV is the slowest model to calibrate and the most stable to hold, so it runs on its own, asynchronous
schedule: the Heston parameters and the conditional expectation are kept until a rule says they no longer describe
the market, while the local vol it multiplies is always today's. Because the leverage is
$\sigma_{\mathrm{LV}}(t, S)/\sqrt{\E[v \mid X]}$, a held table on a fresh local vol still reprices the vanillas
approximately; the error grows with the distance between the market of the table and today's. The tqstudio
Calibration tab has a Heston SLV view (\code{slv\_monitor}, \code{slv\_check} in \code{tqstudio/views/calibration.py})
that compares the latest SLV calibration of each underlying with the latest market and suggests a recalibration
when any rule of \code{SLV\_RULES} fails:
\begin{itemize}[nosep]
  \item age above 7 days;
  \item spot move above 5\%;
  \item surface change above 1 vol point rms (at tenors 3M--3Y and moneyness 80--120\% of the spot, sticky moneyness);
  \item on request, the Monte Carlo fit of the stored model on the latest market above 1 vol point rms
        (\code{slv\_check}, through \code{check\_calibration}).
\end{itemize}
It also shows the drift of the Heston part (closed form) and the ATM 1Y vol then and now. A recalibration today
redoes the whole model (Heston and table); refreshing the table alone in a set (Heston held), and the IMI-style
two-speed Heston with a time-dependent $\varepsilon(t)$, are in \code{ROADMAP.md}.

% ─────────────────────────────────────────────────────────────────────────────
\subsection{Rates calibration}\label{sec:arch-rates-calibration}
\code{calibration/} also holds the rates side, used in the tests and the examples (not in the calibration sets yet):
\begin{itemize}[nosep]
  \item \code{strip\_caplet\_vols(market\_env, generator, quotes)} / \code{cap\_flat\_vols}
        (\code{capletstripping.py}): flat cap vols to caplet vols and back;
  \item \code{calibrate\_gaussian\_short\_rate(market\_env, generator, factors, ...)} (\code{shortrate.py}): Hull--White
        (\code{factors=1}) or G2++ (\code{factors=2}) to caplet and / or swaption vols: model caplets in closed form,
        swaptions batched (approximate or exact), Hull--White $\sigma(t)$ on \code{sigma\_times} (\code{"caplets"}: one
        vol per caplet), \code{fixed} parameters, bounded least squares on vega-scaled residuals with the Jacobian by
        AAD (\code{tf.function}); returns the parameters (with \code{at\_bound}) for a \code{GaussianShortRateModel}.
\end{itemize}
The curve bootstrap is a calibration too, of the curves to the quotes (Section~\ref{sec:arch-ir-data}).
